\section{Sullivan 案时间线：课堂预期与最终程序状态}

课堂用 Uber 2016 年事件讨论 bug bounty、披露和 regulation by enforcement。为了保持教学准确性，必须分四层阅读：Sullivan 在课堂中的个人叙述；政府在起诉与审判中的理论；陪审团和法院已经作出的裁判；本文从中提取的工程控制。任何一层都不能替代另一层。

\subsection{2016--2026 的程序节点}

DOJ 于 2022 年宣布陪审团就妨碍 FTC 程序和 misprision of felony 作出有罪裁决；2023 年法院判处三年 probation。课堂录制时，Sullivan 讨论 \term{nexus}：在 obstruction 语境中，被指行为与官方程序之间需要满足的关联要求，并预期 2024 年最高法院 Fischer 判决可能帮助上诉。第九巡回法院后来维持定罪并发布修订意见；最高法院拒绝 \term{certiorari}（调卷复审请求），即拒绝审理该上诉请求，并不表示最高法院对每个实体争点另行背书。

\lecturefigure{10-current-legal-timeline.png}{课堂中的 pending appeal 已被 2025--2026 年法院程序结果取代。}{DOJ；Ninth Circuit No. 23-927；Supreme Court docket 25-1082；概念重绘。}

\begin{knowledgebox}{读图：把事件、观点与裁判分开}
2016 是底层事件；2022 jury conviction 与 2023 sentencing 是审判结果；2025 Ninth Circuit opinion 是上诉裁判；2026 cert denial 结束本轮最高法院复审请求。课堂关于 nexus、harmless error 或行业影响的陈述属于 2025 年 1 月的诉讼观点。当前讲义可以解释其治理含义，但不能继续写“上诉待决”或把讲者预测当成法院结论。
\end{knowledgebox}

\begin{warningbox}{案件不是通用 VDP 法律测试}
该案涉及特定公司、FTC 调查、事件事实、沟通和记录。不能从一宗刑事案件推出“支付 bounty 都违法”或“签 NDA 都是掩盖”，也不能推出“只要有 VDP 就不会承担责任”。工程上可复用的是分流、升级、证据、披露和共同决策机制；具体法律结论必须回到司法辖区和事实。
\end{warningbox}

\teachervoice{Sullivan 在课堂上把 Fischer 与 nexus 看作上诉的重要机会。保留 teacher voice 的正确方式，是记录“他当时为何这样判断”，同时明确后来法院结果已取代 pending 状态；忠实不是把历史时点冻结成今天的事实。}

\begin{importantbox}{四种陈述标签}
写 case study 时显式区分：\textbf{speaker account}（讲者称）、\textbf{government allegation}（政府主张）、\textbf{adjudicated fact/outcome}（陪审团或法院处理）、\textbf{engineering judgment}（本文建议）。这能减少对人物动机的臆测，也能让读者知道哪一部分可以被系统设计直接采用。
\end{importantbox}

\subsection{本章小结}

这条时间线的教学价值在于纠正旧稿：截至 2026 年 8 月 11 日，上诉不再 pending。更重要的是，组织不能把法律结局当作事件后才需要关注的外部变量；事实记录、权责和披露路径必须在事件发生前设计。

